% Propietario: output/LEY9_SINTESIS_300_SUCESORA_SIN_REGRESIONES_20260904/
% manuscrito/sections/autonomia/03_prolongacion_tpk.tex:354--496.
% Microetapa CG-007. Exposicion autocontenida de dominios y prueba.
\section{Frontière coinductive et stabilisation des prolongements compatibles}
\label{scope:g9-frontera}

L’émission initiale possède une image $\mathcal W_6$ de cardinal $243$ dans l’espace des $729$ mots de six trits. Ses prolongements conservent l’état enrichi : mot, position, direction, feuillet, régime, quotients, retenue, frontière, incidences et registre chronologique. Pour $r\ge1$, notons $\widehat{\mathcal W}_{6r}$ l’ensemble des états admissibles de longueur $6r$ effectivement obtenus par les transitions du TPK, avec leurs préfixes compatibles et leur coordonnée de survie. La projection de mot prend ses valeurs dans $(\mathbb F_3^6)^r$.

La troncature $t_{r+1,r}$ retient le préfixe enrichi de longueur $6r$. La relation de prolongement est
\[
 \mathcal P_r=\{(x,y):x\in\widehat{\mathcal W}_{6r},\
 y\in\widehat{\mathcal W}_{6(r+1)},\ t_{r+1,r}(y)=x\}.
\]
Les relations utilisent les états admissibles produits par le TPK ; l’égalité des préfixes exprime leur compatibilité. Pour $s>r$, écrivons $t_{s,r}=t_{r+1,r}\circ\cdots\circ t_{s,s-1}$. La conservation des champs exige $t_{u,r}=t_{s,r}t_{u,s}$ lorsque $u>s>r$.

\begin{proposition}[Restitution de l’antécédent par troncature]
Si $p_r$ est une section de prolongement dont le graphe est contenu dans $\mathcal P_r$, alors $t_{r+1,r}p_r=I$. Pour une suite de telles sections, la composition de leurs troncatures récupère l’état antécédent complet.
\end{proposition}
\begin{proof}
L’inclusion du graphe signifie, pour chaque $x$ du domaine de la section, $t_{r+1,r}(p_r(x))=x$. Lors de la composition de deux sections, le couple de plus grande profondeur s’annule d’abord, puis le précédent. La récurrence établit la formule pour tout nombre fini d’étapes. L’égalité porte sur les états enrichis et conserve donc tous leurs champs et leurs relations de préfixe.
\end{proof}

\subsection{Candidats à la première frontière et relation d’ouverture}

Soit $\widetilde{\mathcal R}_{36}$ l’ensemble des états admissibles de profondeur $36$ dont le registre de frontière et la mémoire de bord sont définis. Pour $x\in\widetilde{\mathcal R}_{36}$, l’ensemble $\operatorname{Cand}_6(x)$ est constitué des états $y$ de profondeur $42$ obtenus par une trajectoire typée
\[
 x=x_{36}\xrightarrow{\mathcal U_{36}}x_{37}
 \xrightarrow{\mathcal U_{37}}\cdots
 \xrightarrow{\mathcal U_{41}}x_{42}=y,
 \qquad \tau_{42,36}(y)=x.
\]
La queue visible de chaque candidat comporte six trits ; sa projection de mot possède donc au plus $729$ valeurs. Le cardinal de cette projection se distingue de celui des états à registres complets.

Le prédicat $g_{9,x}(y)$ exige la compatibilité de tous les préfixes et l’existence de continuations admissibles à profondeur arbitraire. La frontière effective et son ouverture sont définies par
\begin{align*}
 \mathcal R_{36}
 &=\{x\in\widetilde{\mathcal R}_{36}:
      \exists y\in\operatorname{Cand}_6(x),\ g_{9,x}(y)=1\},\\
 G_9&=\{(x,y):x\in\mathcal R_{36},\
                y\in\operatorname{Cand}_6(x),\ g_{9,x}(y)=1\}.
\end{align*}
La construction conserve le carré de troncature $\tau_{42,36}(y)=x$ et donne
\[
 G_9\subseteq\mathcal U_{41}\circ\cdots\circ\mathcal U_{36}
 \subseteq\widehat{\mathcal X}_{36}\times\widehat{\mathcal X}_{42}.
\]
Chaque élément de $\mathcal R_{36}$ possède par définition un candidat survivant. La relation $G_9$ ouvre un bloc de six trits ; l’holonomie $\Gamma_9$ compose neuf transitions élémentaires. Le premier objet est une relation entre niveaux de profondeur, et le second est le retour du calendrier avec son incrément de mémoire.

\subsection{Existence d’une trajectoire complète}

\begin{theorem}[Stabilisation avec conservation de la mémoire]
Supposons que les préfixes admissibles proviennent d’un ensemble fini non vide de conditions initiales et forment un système à branchement fini, avec des états à toute profondeur, dont les troncatures satisfont la compatibilité précédente et conservent les champs de l’état. Supposons également la naturalité du transport nonadique par rapport à la troncature et les identités de retour $g_{k+9}=g_k$, $m_{k+9}=m_k+1$ sur les trajectoires admises. Il existe alors une trajectoire infinie compatible. Le long de celle-ci, chaque retour conserve la phase et augmente la mémoire entière d’une unité.
\end{theorem}
\begin{proof}
Chaque préfixe est organisé selon sa longueur et relié à ses prolongements immédiats. S’il existe plusieurs états initiaux, on ajoute une racine formelle possédant cet ensemble fini de successeurs. La non-vacuité à toute profondeur fournit des extensions de longueur arbitraire à partir de cette racine.

Un sommet admettant des extensions arbitrairement longues possède au moins un successeur ayant la même propriété. En effet, si chacun de ses successeurs, en nombre fini, avait une borne finie de profondeur, le maximum de ces bornes, augmenté d’un, bornerait toutes les extensions du sommet. Cela contredit la propriété initiale. On choisit successivement un successeur à extension illimitée ; l’ordre fini des mots et des registres permet de fixer le choix lorsque ce codage est disponible. L’itération produit une chaîne ayant un préfixe à chaque niveau.

La relation de prolongement garantit que chaque troncature immédiate restitue l’antécédent. Sa compatibilité par composition démontre la même égalité entre deux profondeurs quelconques. La naturalité permet d’appliquer l’identité de retour à tous les préfixes où elle est définie. Ses deux composantes fournissent respectivement la phase conservée et l’incrément unitaire de mémoire. Après $s$ tours, $g_{k+9s}=g_k$ et $m_{k+9s}=m_k+s$, par récurrence.
\end{proof}

La démonstration établit l’existence à partir du branchement fini, du prolongement à toute profondeur et de la compatibilité. Le prédicat de survie exprime cette propriété de la trajectoire complète. L’individualisation d’une section régionale concrète utilise en outre le lecteur et la règle fonctionnelle de cette région, développés dans sa propre chaîne. Ces deux résultats conservent leurs domaines respectifs et se composent lorsque la section sélectionnée satisfait les relations de prolongement précédentes.
